%% Chapter 5 — RFID Card Architecture
\chapter{RFID Card Architecture}
\label{ch:rfid}

\section{MIFARE Classic 1K Overview}

MIFARE Classic 1K cards have 1024 bytes of non-volatile EEPROM organised as:
\begin{itemize}[leftmargin=1.5em]
    \item 16 sectors (sectors 0--15)
    \item 4 blocks per sector (each block = 16 bytes)
    \item Block 0 of Sector 0 is the manufacturer block (UID + config; read-only)
    \item Each sector has a trailer block (last block in sector) containing two keys (Key A and Key B) and access bits
    \item Data blocks: 3 per sector (excluding trailer)
    \item Sectors are authenticated independently using their own keys
\end{itemize}

\section{V2 Card Schema}

MMS V2 uses \textbf{Sector 1} (blocks 4--7). Sector 0 is left untouched (manufacturer data and default keys). Using Sector 1 means V2 cards can coexist with other applications on Sector 0 if needed.

\subsection{Sector 1 Layout}

\begin{table}[H]
\centering
\caption{MMS V2 RFID Card Schema --- Sector 1}
\label{tab:card_schema}
\begin{tabular}{C{1.5cm}C{1.5cm}L{10.5cm}}
\toprule
\textbf{Block} & \textbf{Bytes} & \textbf{Content} \\
\midrule
4 & 0--8   & Roll number as ASCII string, zero-padded to 9 characters \\
  & 9      & Schema version byte: \texttt{0x02} (V2) \\
  & 10--15 & Reserved, all \texttt{0x00} \\
\midrule
5 & 0--7   & PIN hash: $\text{HMAC-SHA256}(\text{master\_key} \| \text{uid\_hex},\, \text{pin\_string})[0..7]$ \\
  & 8--11  & Machine permission bitmask (uint32, little-endian) \\
  & 12--15 & Reserved, all \texttt{0x00} \\
\midrule
6 & 0--15  & Reserved for future use (expiry date, academic year) \\
\midrule
7 & 0--5   & Key A = derived sector key (write-only; never returned by MFRC522 on read) \\
  & 6--9   & Access bits \\
  & 10--15 & Key B = \texttt{0x00} (not used) \\
\bottomrule
\end{tabular}
\end{table}

\subsection{Schema Version Byte}

The schema version byte at Block 4, byte 9 is \texttt{0x02} for all V2 cards. If a card is read and this byte is not \texttt{0x02}, the firmware rejects it with \texttt{ACCESS\_DENIED} and logs \texttt{"schema version mismatch"}. This prevents:
\begin{itemize}[leftmargin=1.5em]
    \item V1 cards (written to Block 1+2 with no schema byte) from being accepted
    \item Blank or foreign cards from accidentally gaining access
    \item Future V3 cards (with schema byte \texttt{0x03}) from being accepted by V2 nodes
\end{itemize}

\section{Sector Key Derivation}

Each card's Sector 1 is protected by a unique 6-byte key derived from the master key and the card's UID:

\[
\text{sector\_key}[6] = \text{HMAC-SHA256}\!\left(\text{master\_key},\, \text{uid\_as\_uppercase\_hex\_string}\right)[0..5]
\]

For a card with UID bytes \texttt{[0xA1, 0xB2, 0xC3, 0xD4]}, the UID string is \texttt{"A1B2C3D4"}.

This derivation is performed by:
\begin{itemize}[leftmargin=1.5em]
    \item The kiosk (\texttt{station\_writer.ino}) when writing the card, to set the sector key
    \item Every access node on every card read, to authenticate Sector 1 before reading blocks
\end{itemize}

The master key is identical across all nodes and the kiosk --- they all share the same 32-byte secret stored in \texttt{secrets.h} / \texttt{secrets.py} / \texttt{secrets\_station.h}.

\begin{notebox}
HMAC-SHA256 is implemented using \texttt{mbedtls\_md\_hmac()} from the ESP32's built-in mbedTLS library. No additional library is required. The same function is called from the kiosk (C++ Arduino) and the RPi bridge (Python \texttt{hmac} module with sha256 digest).
\end{notebox}

\section{PIN Hash}

The PIN hash stored in Block 5 bytes 0--7 is:
\[
\text{pin\_hash}[8] = \text{HMAC-SHA256}\!\left(\text{master\_key} \| \text{uid\_hex},\, \text{pin\_string}\right)[0..7]
\]

Where \texttt{pin\_string} is the 4-character ASCII PIN (e.g., ``1234'').

The PIN hash is computed by:
\begin{itemize}[leftmargin=1.5em]
    \item The kiosk when issuing the card
    \item The web app's \texttt{hashPin()} function when the admin sets a PIN
    \item Access nodes do NOT verify PIN in V2 (Phase 4)
\end{itemize}

\begin{warnbox}
\textbf{In V2, the PIN is not verified at the machine.} The PIN hash is stored on the card for future use (Phase 4 --- rotary encoder PIN entry). In V2, card tap + permission bit check is sufficient for access.
\end{warnbox}

\section{Machine Permission Bitmask}

Block 5 bytes 8--11 contain a 32-bit little-endian unsigned integer. Each bit corresponds to one machine:

\begin{equation}
\text{allowed} = (\text{perm\_bitmask} \gg (\text{machine\_id} - 1)) \;\&\; 1
\end{equation}

\begin{table}[H]
\centering
\caption{Machine Permission Bitmask Examples}
\begin{tabular}{L{3cm}C{2cm}C{3cm}L{5.5cm}}
\toprule
\textbf{Access to machines} & \textbf{Decimal} & \textbf{Hex} & \textbf{Binary (LSB first)} \\
\midrule
None (no access) & 0 & \texttt{0x00000000} & \texttt{00000000...} \\
Machine 1 only & 1 & \texttt{0x00000001} & \texttt{10000000...} (bit 0 set) \\
Machine 3 only & 4 & \texttt{0x00000004} & \texttt{00100000...} (bit 2 set) \\
Machines 1 and 3 & 5 & \texttt{0x00000005} & \texttt{10100000...} (bits 0+2 set) \\
All 5 machines & 31 & \texttt{0x0000001F} & \texttt{11111000...} (bits 0--4 set) \\
Full access (all 32) & 4294967295 & \texttt{0xFFFFFFFF} & all bits set \\
\bottomrule
\end{tabular}
\end{table}

The bitmask supports up to 32 machines. Currently 5 are deployed (bits 0--4). Bits 5--31 are reserved for future machines.

\section{Card Write Process}

When the kiosk writes a card:

\begin{enumerate}[leftmargin=1.5em]
    \item Admin selects student and machine permissions in web UI
    \item Flask app computes bitmask and PIN hash
    \item Flask app sends JSON to station\_writer via serial: \\
          \texttt{\{"roll":"220002123","perm\_mask":5,"pin\_hash":"aabb...","uid":"A1B2C3D4"\}}
    \item station\_writer.ino receives JSON (ArduinoJson), places card on MFRC522
    \item Derives sector key from UID
    \item Authenticates Sector 1 with default key (\texttt{0xFFFFFFFFFFFF} for fresh cards) or derived key (for re-issue)
    \item Writes Block 4: roll bytes 0--8 + schema byte \texttt{0x02} + zeros 10--15
    \item Writes Block 5: pin\_hash bytes 0--7 + perm\_mask LE bytes 8--11 + zeros 12--15
    \item Writes Sector 1 trailer (Block 7): Key A = derived sector key, Key B = zeros
    \item Reads back Block 4 and verifies schema byte = \texttt{0x02}
    \item Reports \texttt{WRITE\_SUCCESS} or \texttt{WRITE\_FAILED} via serial to Flask app
\end{enumerate}

\section{Card Read Process}

When a node reads a card during an access attempt:

\begin{enumerate}[leftmargin=1.5em]
    \item HC89 signals card present
    \item \texttt{MFRC522::PICC\_IsNewCardPresent()} confirms a card in field
    \item \texttt{MFRC522::PICC\_ReadCardSerial()} reads the 4-byte UID
    \item Sector key is derived from UID + master key
    \item \texttt{MFRC522::PCD\_Authenticate()} authenticates Sector 1 with derived key
    \item Block 4 is read; schema byte checked (must be \texttt{0x02})
    \item Block 5 is read; perm\_mask extracted from bytes 8--11 (LE decode)
    \item UID is checked against LittleFS revocation list
    \item Permission bit checked for this machine\_id
    \item Result: GRANTED or DENIED with reason
\end{enumerate}
